% Generated draft: requires mathematical and source review.
\documentclass[11pt]{article}
\usepackage[margin=25mm]{geometry}
\usepackage{fontspec}
\setmainfont{Latin Modern Roman}
\usepackage{amsmath,amssymb,mathtools,unicode-math}
\setmathfont{Latin Modern Math}
\usepackage{enumitem,needspace,xurl,hyperref}
\setlength{\parindent}{0pt}
\setlength{\parskip}{5pt}
\setlength{\emergencystretch}{3em}
\newcommand{\N}{\mathbb{N}}
\newcommand{\Z}{\mathbb{Z}}
\newcommand{\Q}{\mathbb{Q}}
\newcommand{\R}{\mathbb{R}}
\newcommand{\C}{\mathbb{C}}
\newcommand{\F}{\mathbb{F}}
\newcommand{\A}{\mathbb{A}}
\newcommand{\PP}{\mathbb{P}}
\newcommand{\E}{\mathbb{E}}
\newcommand{\eps}{\varepsilon}
\newcommand{\dd}{\,\mathrm{d}}
\begin{document}
\section*{UnsolvedMath ID: 6800002}
\textbf{Source problem code:} catalogue entry
\subsection{Definitions and mathematical statement}
Write \(\operatorname{inj}(M,g)=\inf_{x\in M}\operatorname{inj}_g(x)\). Interpret the source's Ricci lower bound as the quadratic-form inequality \(\operatorname{Ric}_g\ge k g\). The proposed assertion is
\[
 \forall\ell,k>0\ \exists C,L,K>0\ \forall(M,g):
\]
if \((M,g)\) is complete with \(\operatorname{inj}(M,g)\ge\ell\) and \(\operatorname{Ric}_g\ge k g\), then there is a Riemannian metric \(h\) on the same manifold such that
\[
 \operatorname{inj}(M,h)\ge L,\qquad |\operatorname{Ric}_h|_h\le K,
 \qquad C^{-1}g\le h\le Cg.
\]
The metric comparison means \(C^{-1}g_x(v,v)\le h_x(v,v)\le Cg_x(v,v)\) for every tangent vector. The constants are uniform over the admitted manifolds, and may be enlarged so \(C\ge1\). The positive sign in \(k>0\) and \(\operatorname{Ric}_g\ge k g\) is retained exactly; no negative Ricci lower bound is substituted. The conclusion does not require \(h\) to retain positive Ricci curvature.

No dimension parameter appears in the supplied statement. Its literal quantifier pattern permits dependence only on \(\ell,k\); whether a dimension was fixed in surrounding context, allowing constants also to depend on it, requires confirmation. The source also leaves \(|\operatorname{Ric}_h|\) as an unspecified tensor norm. Operator and Hilbert--Schmidt conventions give different dimension dependence, so one cannot silently interchange them in a dimension-uniform reading.

The supplied summary reports related metric-perturbation and operator-stability results by Bandara and collaborators, but says those use a nearby controlled metric as a hypothesis rather than prove this existence statement. It reports no verified solution of the exact question. Those are source assessments, not an independently checked current status.
\subsection{Short English statement}
Can a complete metric with uniformly positive Ricci curvature and positive injectivity radius be replaced by a uniformly comparable metric with two-sided bounded Ricci curvature and a uniform positive injectivity radius? The constants must have the source's stated dependence; dimension and tensor-norm conventions need confirmation.
\subsection{Sources}
\begin{itemize}
\item Local registry: \url{data/unsolvedmath/problems.json}, permanent UnsolvedMath ID \(6800002\), source code catalogue entry; supplied statement, background, and research summary, updated 2026-08-06. Proposed by Lashi Bandara, Universit\"at Potsdam, Germany, 2018.
\item Morgan and Pansu, \emph{A List of Open Problems in Differential Geometry} (2018), Question environment 2, ``Riemannian manifolds with curvature bounds'': \url{https://www.imo.universite-paris-saclay.fr/~pansu/problems_MTDG.tex}. Supplied source-TeX extraction accessed 2026-07-29.
\item Bandara--McIntosh--Ros\'en, ``Riesz continuity of the Atiyah--Singer Dirac operator under perturbations of the metric'' (2017), as supplied by the summary.
\item Bandara--McIntosh, ``Rough metrics on manifolds and quadratic estimates'', arXiv:1402.2030, as supplied by the summary. No exact resolving theorem is included.
\end{itemize}
\end{document}
